% Suplemento de trabajo. No sustituye registro_k.tex ni cierra por decreto E108.
% Procedencia y alcance: INFORME_RECUPERACION.md, misma carpeta.
\subsection{Évaluation incidencielle antérieure au registre dodécaphasique}
\label{subsec:registro-evaluacion-incidencial}

La récupération d'un registre à partir de ses différences cycliques et de sa
charge totale est une question distincte de l'évaluation de ces coordonnées
sur le livre d'incidences. Le premier problème se résout au moyen de
l'extracteur transversal de la proposition~\ref{prop:k-transversal}.
Pour le second, une présentation historique du
registre fournit trois dénombrements entiers dans chaque fenêtre. Il est utile
d'expliciter leur composition et de déterminer exactement l'information que
requiert leur évaluation.

Soit \(\mathscr L\) un livre de visites et de traces orientées préalablement
généré. Chaque visite conserve la route, la position APP, le feuillet,
l'instant, la multiplicité et l'incidence de frontière. L'évaluation
arithmétique produit l'entier et sa décomposition
\[
 n_t=r_t+9q_t,\qquad 1\leq r_t\leq9.
\]
L'entier \(n_t\) se calcule à partir du feuillet de somme ou de produit ; son
résidu n'est pas reçu comme une valeur indépendante. La distinction entre
incidence de couronne et incidence intérieure demeure explicite pour les
visites de résidu \(9\).

Dans la fenêtre \(E_m\), la présentation par dénombrements considère
\begin{align}
 A_m&=\sum_{v\in E_m}w(v)
             \mathbf1_{\{1,3,5,7\}}(r(v)),\\
 C_m&=\sum_{j\geq0}\bigl(\nu^-_{120,m}(j)-\nu^+_{120,m}(j)\bigr),\\
 V_m&=\sum_{v\in E_m}w(v)\mathbf1_{\{9\}}(r(v))\epsilon_\partial(v).
\end{align}
Ici, \(w(v)\) est la multiplicité incidencielle et
\(\epsilon_\partial(v)=1\) sur la couronne, \(-1\) à l'intérieur.
Les traces comptées par \(\nu^\pm\) doivent être énumérées au moyen de la règle
correspondante du livre. La somme entière exige un support fini ou une
construction alternative qui détermine sa signification. Le nombre
d'incidences chronologiques et la longueur réduite \(120(1+3j)\) sont de
types distincts ; leur relation exige l'unité de longueur du lecteur.

On définit \([x]_{10}\in\{0,\ldots,9\}\) comme le représentant
euclidien, et l'on conserve les quotients
\[
 x=[x]_{10}+10\lfloor x/10\rfloor.
\]
Le bloc incidenciel est
\begin{equation}
 z_m=100[A_m]_{10}+10[C_m]_{10}+[V_m]_{10}.
 \label{eq:registro-bloque-incidencial}
\end{equation}
Par composition, on obtient l'application
\begin{equation}
 \mathcal E_{\rm cnt}(\mathscr L)=
 \left((z_m-z_{m+3})_{m=1}^{12},
       (z_m-z_{m+4})_{m=1}^{12},\sum_{m=1}^{12}z_m\right).
\label{eq:registro-composicion-incidencial}
\end{equation}
Les indices se lisent cycliquement modulo \(12\). Dans cette partie,
\((D_jz)_m=z_m-z_{m+j}\), pour \(j=3,4\), et
\(Q(z)=\sum_{m=1}^{12}z_m\) ; par conséquent, la composition précédente est
\((D_3z,D_4z,Q(z))\). Cette convention est conservée dans
\S\ref{subsec:k-transversal}, où est démontrée la reconstruction
entière complète.
Cette formule produit ses coordonnées à partir des dénombrements. Elle n’utilise
ni \(K\), ni \(U\), ni \(\alpha\), ni les coordonnées transversales attendues
comme arguments. Son identification à
\(\mathcal E_{108}^{90,120}\) nécessite de la composer avec la famille et
les incidences effectives de l’état terminal ; elle ne se déduit pas uniquement
de l’inversibilité du système transversal.

\begin{proposition}[Information arithmétique suffisante et nécessaire]
\label{prop:registro-36-residuos}
Soient \(N,N'\in\mathbb Z^{12\times3}\), ordonnés selon les dénombrements
\((A,C,V)\). Pour l'application définie par
\eqref{eq:registro-bloque-incidencial}--\eqref{eq:registro-composicion-incidencial},
\[
 \mathcal E_{\rm cnt}(N)=\mathcal E_{\rm cnt}(N')
 \quad\Longleftrightarrow\quad
 N-N'\in10\mathbb Z^{36}.
\]
En particulier, les \(36\) résidus sectoriels déterminent exactement
la sortie à \(25\) coordonnées. La conservation des quotients et
de la provenance constitue une information supplémentaire du registre.
\end{proposition}

\begin{proof}
L'égalité modulo \(10\) produit les mêmes blocs et, par conséquent,
les mêmes coordonnées transversales. Réciproquement, si les coordonnées
transversales coïncident, \(D_3(z-z')=D_4(z-z')=0\). Les pas
\(3\) et \(4\) engendrent le cycle de longueur \(12\), de sorte que
\(z-z'\) est constant. L'égalité des charges totales annule cette
constante. Enfin, la représentation décimale d'un entier compris entre
\(0\) et \(999\) possède trois chiffres uniques. On récupère ainsi les
trois résidus de chaque fenêtre. L'assertion est une identité algébrique
pour tous \(N,N'\), et non une inférence obtenue par des essais finis.
\end{proof}

\subsection{Conjugaison des dénombrements et termes de frontière}
\label{subsec:registro-conjugacion-incidencial}

Soit \(\rho(m)=13-m\). Considérons une conjugaison d'incidences
qui inverse l'ordre des fenêtres, conserve la classe arithmétique et la
classe de frontière des visites et échange les canaux des traces.
Alors
\[
 (A_m^\dagger,C_m^\dagger,V_m^\dagger)
 =(A_{\rho(m)},-C_{\rho(m)},V_{\rho(m)}).
\]
Si \(c_m=[C_m]_{10}\), la réduction décimale produit l'identité
\begin{equation}
 z_m^\dagger=z_{\rho(m)}
       +100\mathbf1_{\{c_{\rho(m)}\ne0\}}-20c_{\rho(m)}.
 \label{eq:registro-conjugacion-decimal}
\end{equation}
En effet, \([-C]_{10}=0\) lorsque \([C]_{10}=0\), et
\([-C]_{10}=10-[C]_{10}\) dans le cas contraire. La formule conserve
exactement cette différence. La négation du canal n'équivaut pas à nier
le bloc décimal complet.

L'application de cette identité à une involution concrète du TPK exige
de vérifier son action sur les incidences. Pour un lecteur d'occupations
évalué avant l'avancement, l'inversion d'une route
\(\gamma=(x_0,\ldots,x_n)\) satisfait
\[
 \sum_{k=0}^{n-1}f(x_{n-k})
 =\sum_{k=0}^{n-1}f(x_k)+f(x_n)-f(x_0).
\]
Les termes aux extrémités sont incorporés avant la réduction modulo
\(10\). Ils s'annulent sur les routes fermées ; dans les fenêtres ouvertes, ils peuvent
modifier leurs résidus. Cette correction, déjà présente dans le développement
bidirectionnel du transport, spécifie ce que doit conserver le raccordement
entre la route et les dénombrements.

\begin{proposition}[Nécessité de l'information non périodique du registre]
Supposons que l’observable d’une famille fixe de chemins satisfasse
\(o(t+54)=o(t)\), que la multiplicité soit conservée sous ce retour
et que le même lecteur local agisse dans les fenêtres de \(9\) instants.
Alors ses dénombrements et blocs satisfont
\[
 (A,C,V)_{m+6}=(A,C,V)_m,\qquad z_{m+6}=z_m.
\]
\end{proposition}

\begin{proof}
La translation \(t\mapsto t+54\) envoie \(E_m\) bijectivement
sur \(E_{m+6}\), en conservant chaque incidence comptée et sa multiplicité.
L'égalité se transmet à la somme et à la réduction décimale.
\end{proof}

Une présentation à \(12\) blocs distincts exige donc que
le lecteur conserve une information qui ne se factorise pas par cet observable
de période \(54\) : sélection d'incidences, orientation effective,
extrémités, multiplicité ou mémoire. La proposition délimite une réduction
inadmissible de l'état ; elle n'affirme pas que la dynamique complète possède une
période \(54\) et ne détermine pas à elle seule son lecteur terminal.
